\section{}

\paragraph{982年}\eventanchor{XIA-0982-EVENT-049}{XIA-0982-EVENT-049}　李继迁脱离宋朝羁縻控制，在西北边地集结力量。此项以《宋史》卷485〈夏国传〉；《辽史》《金史》相关本纪；黑水城西夏文文书与考古报告 为基本来源线索；编年与官修材料能够确认年序和制度用语，但对地方执行、人数、动机或社会后果的判断仍须同文书、碑刻、考古及对方叙事互校。

\paragraph{982年}\eventanchor{XIA-0982-REGNAL-YEAR}{XIA-0982-REGNAL-YEAR}　党项李氏以李继迁在位第1年作为诏令、赋役与外交文书的纪年框架。此项以《宋史》卷485〈夏国传〉；《辽史》《金史》相关本纪；黑水城西夏文文书与考古报告 为基本来源线索；编年与官修材料能够确认年序和制度用语，但对地方执行、人数、动机或社会后果的判断仍须同文书、碑刻、考古及对方叙事互校。

\paragraph{983年}\eventanchor{XIA-0983-REGNAL-YEAR}{XIA-0983-REGNAL-YEAR}　党项李氏以李继迁在位第2年作为诏令、赋役与外交文书的纪年框架。此项以《宋史》卷485〈夏国传〉；《辽史》《金史》相关本纪；黑水城西夏文文书与考古报告 为基本来源线索；编年与官修材料能够确认年序和制度用语，但对地方执行、人数、动机或社会后果的判断仍须同文书、碑刻、考古及对方叙事互校。

\paragraph{984年}\eventanchor{XIA-0984-REGNAL-YEAR}{XIA-0984-REGNAL-YEAR}　党项李氏以李继迁在位第3年作为诏令、赋役与外交文书的纪年框架。此项以《宋史》卷485〈夏国传〉；《辽史》《金史》相关本纪；黑水城西夏文文书与考古报告 为基本来源线索；编年与官修材料能够确认年序和制度用语，但对地方执行、人数、动机或社会后果的判断仍须同文书、碑刻、考古及对方叙事互校。

\paragraph{985年}\eventanchor{XIA-0985-REGNAL-YEAR}{XIA-0985-REGNAL-YEAR}　党项李氏以李继迁在位第4年作为诏令、赋役与外交文书的纪年框架。此项以《宋史》卷485〈夏国传〉；《辽史》《金史》相关本纪；黑水城西夏文文书与考古报告 为基本来源线索；编年与官修材料能够确认年序和制度用语，但对地方执行、人数、动机或社会后果的判断仍须同文书、碑刻、考古及对方叙事互校。

\paragraph{986年}\eventanchor{XIA-0986-REGNAL-YEAR}{XIA-0986-REGNAL-YEAR}　党项李氏以李继迁在位第5年作为诏令、赋役与外交文书的纪年框架。此项以《宋史》卷485〈夏国传〉；《辽史》《金史》相关本纪；黑水城西夏文文书与考古报告 为基本来源线索；编年与官修材料能够确认年序和制度用语，但对地方执行、人数、动机或社会后果的判断仍须同文书、碑刻、考古及对方叙事互校。

\paragraph{987年}\eventanchor{XIA-0987-REGNAL-YEAR}{XIA-0987-REGNAL-YEAR}　党项李氏以李继迁在位第6年作为诏令、赋役与外交文书的纪年框架。此项以《宋史》卷485〈夏国传〉；《辽史》《金史》相关本纪；黑水城西夏文文书与考古报告 为基本来源线索；编年与官修材料能够确认年序和制度用语，但对地方执行、人数、动机或社会后果的判断仍须同文书、碑刻、考古及对方叙事互校。

\paragraph{988年}\eventanchor{XIA-0988-REGNAL-YEAR}{XIA-0988-REGNAL-YEAR}　党项李氏以李继迁在位第7年作为诏令、赋役与外交文书的纪年框架。此项以《宋史》卷485〈夏国传〉；《辽史》《金史》相关本纪；黑水城西夏文文书与考古报告 为基本来源线索；编年与官修材料能够确认年序和制度用语，但对地方执行、人数、动机或社会后果的判断仍须同文书、碑刻、考古及对方叙事互校。

\paragraph{989年}\eventanchor{XIA-0989-REGNAL-YEAR}{XIA-0989-REGNAL-YEAR}　党项李氏以李继迁在位第8年作为诏令、赋役与外交文书的纪年框架。此项以《宋史》卷485〈夏国传〉；《辽史》《金史》相关本纪；黑水城西夏文文书与考古报告 为基本来源线索；编年与官修材料能够确认年序和制度用语，但对地方执行、人数、动机或社会后果的判断仍须同文书、碑刻、考古及对方叙事互校。

\paragraph{990年}\eventanchor{XIA-0990-EVENT-050}{XIA-0990-EVENT-050}　李继迁在战与和之间接受宋廷封授并保持独立军事行动。此项以《宋史》卷485〈夏国传〉；《辽史》《金史》相关本纪；黑水城西夏文文书与考古报告 为基本来源线索；编年与官修材料能够确认年序和制度用语，但对地方执行、人数、动机或社会后果的判断仍须同文书、碑刻、考古及对方叙事互校。

\paragraph{990年}\eventanchor{XIA-0990-REGNAL-YEAR}{XIA-0990-REGNAL-YEAR}　党项李氏以李继迁在位第9年作为诏令、赋役与外交文书的纪年框架。此项以《宋史》卷485〈夏国传〉；《辽史》《金史》相关本纪；黑水城西夏文文书与考古报告 为基本来源线索；编年与官修材料能够确认年序和制度用语，但对地方执行、人数、动机或社会后果的判断仍须同文书、碑刻、考古及对方叙事互校。

\paragraph{991年}\eventanchor{XIA-0991-REGNAL-YEAR}{XIA-0991-REGNAL-YEAR}　党项李氏以李继迁在位第10年作为诏令、赋役与外交文书的纪年框架。此项以《宋史》卷485〈夏国传〉；《辽史》《金史》相关本纪；黑水城西夏文文书与考古报告 为基本来源线索；编年与官修材料能够确认年序和制度用语，但对地方执行、人数、动机或社会后果的判断仍须同文书、碑刻、考古及对方叙事互校。

\paragraph{992年}\eventanchor{XIA-0992-REGNAL-YEAR}{XIA-0992-REGNAL-YEAR}　党项李氏以李继迁在位第11年作为诏令、赋役与外交文书的纪年框架。此项以《宋史》卷485〈夏国传〉；《辽史》《金史》相关本纪；黑水城西夏文文书与考古报告 为基本来源线索；编年与官修材料能够确认年序和制度用语，但对地方执行、人数、动机或社会后果的判断仍须同文书、碑刻、考古及对方叙事互校。

\paragraph{993年}\eventanchor{XIA-0993-REGNAL-YEAR}{XIA-0993-REGNAL-YEAR}　党项李氏以李继迁在位第12年作为诏令、赋役与外交文书的纪年框架。此项以《宋史》卷485〈夏国传〉；《辽史》《金史》相关本纪；黑水城西夏文文书与考古报告 为基本来源线索；编年与官修材料能够确认年序和制度用语，但对地方执行、人数、动机或社会后果的判断仍须同文书、碑刻、考古及对方叙事互校。

\paragraph{994年}\eventanchor{XIA-0994-REGNAL-YEAR}{XIA-0994-REGNAL-YEAR}　党项李氏以李继迁在位第13年作为诏令、赋役与外交文书的纪年框架。此项以《宋史》卷485〈夏国传〉；《辽史》《金史》相关本纪；黑水城西夏文文书与考古报告 为基本来源线索；编年与官修材料能够确认年序和制度用语，但对地方执行、人数、动机或社会后果的判断仍须同文书、碑刻、考古及对方叙事互校。

\paragraph{995年}\eventanchor{XIA-0995-REGNAL-YEAR}{XIA-0995-REGNAL-YEAR}　党项李氏以李继迁在位第14年作为诏令、赋役与外交文书的纪年框架。此项以《宋史》卷485〈夏国传〉；《辽史》《金史》相关本纪；黑水城西夏文文书与考古报告 为基本来源线索；编年与官修材料能够确认年序和制度用语，但对地方执行、人数、动机或社会后果的判断仍须同文书、碑刻、考古及对方叙事互校。

\paragraph{996年}\eventanchor{XIA-0996-REGNAL-YEAR}{XIA-0996-REGNAL-YEAR}　党项李氏以李继迁在位第15年作为诏令、赋役与外交文书的纪年框架。此项以《宋史》卷485〈夏国传〉；《辽史》《金史》相关本纪；黑水城西夏文文书与考古报告 为基本来源线索；编年与官修材料能够确认年序和制度用语，但对地方执行、人数、动机或社会后果的判断仍须同文书、碑刻、考古及对方叙事互校。

\paragraph{997年}\eventanchor{XIA-0997-REGNAL-YEAR}{XIA-0997-REGNAL-YEAR}　党项李氏以李继迁在位第16年作为诏令、赋役与外交文书的纪年框架。此项以《宋史》卷485〈夏国传〉；《辽史》《金史》相关本纪；黑水城西夏文文书与考古报告 为基本来源线索；编年与官修材料能够确认年序和制度用语，但对地方执行、人数、动机或社会后果的判断仍须同文书、碑刻、考古及对方叙事互校。

\paragraph{998年}\eventanchor{XIA-0998-REGNAL-YEAR}{XIA-0998-REGNAL-YEAR}　党项李氏以李继迁在位第17年作为诏令、赋役与外交文书的纪年框架。此项以《宋史》卷485〈夏国传〉；《辽史》《金史》相关本纪；黑水城西夏文文书与考古报告 为基本来源线索；编年与官修材料能够确认年序和制度用语，但对地方执行、人数、动机或社会后果的判断仍须同文书、碑刻、考古及对方叙事互校。

\paragraph{999年}\eventanchor{XIA-0999-REGNAL-YEAR}{XIA-0999-REGNAL-YEAR}　党项李氏以李继迁在位第18年作为诏令、赋役与外交文书的纪年框架。此项以《宋史》卷485〈夏国传〉；《辽史》《金史》相关本纪；黑水城西夏文文书与考古报告 为基本来源线索；编年与官修材料能够确认年序和制度用语，但对地方执行、人数、动机或社会后果的判断仍须同文书、碑刻、考古及对方叙事互校。

\paragraph{1000年}\eventanchor{XIA-1000-REGNAL-YEAR}{XIA-1000-REGNAL-YEAR}　党项李氏以李继迁在位第19年作为诏令、赋役与外交文书的纪年框架。此项以《宋史》卷485〈夏国传〉；《辽史》《金史》相关本纪；黑水城西夏文文书与考古报告 为基本来源线索；编年与官修材料能够确认年序和制度用语，但对地方执行、人数、动机或社会后果的判断仍须同文书、碑刻、考古及对方叙事互校。

\paragraph{1001年}\eventanchor{XIA-1001-REGNAL-YEAR}{XIA-1001-REGNAL-YEAR}　党项李氏以李继迁在位第20年作为诏令、赋役与外交文书的纪年框架。此项以《宋史》卷485〈夏国传〉；《辽史》《金史》相关本纪；黑水城西夏文文书与考古报告 为基本来源线索；编年与官修材料能够确认年序和制度用语，但对地方执行、人数、动机或社会后果的判断仍须同文书、碑刻、考古及对方叙事互校。

\paragraph{1002年}\eventanchor{XIA-1002-REGNAL-YEAR}{XIA-1002-REGNAL-YEAR}　党项李氏以李继迁在位第21年作为诏令、赋役与外交文书的纪年框架。此项以《宋史》卷485〈夏国传〉；《辽史》《金史》相关本纪；黑水城西夏文文书与考古报告 为基本来源线索；编年与官修材料能够确认年序和制度用语，但对地方执行、人数、动机或社会后果的判断仍须同文书、碑刻、考古及对方叙事互校。

\paragraph{1003年}\eventanchor{XIA-1003-REGNAL-YEAR}{XIA-1003-REGNAL-YEAR}　党项李氏以李继迁在位第22年作为诏令、赋役与外交文书的纪年框架。此项以《宋史》卷485〈夏国传〉；《辽史》《金史》相关本纪；黑水城西夏文文书与考古报告 为基本来源线索；编年与官修材料能够确认年序和制度用语，但对地方执行、人数、动机或社会后果的判断仍须同文书、碑刻、考古及对方叙事互校。

\paragraph{1004年}\eventanchor{XIA-1004-EVENT-051}{XIA-1004-EVENT-051}　李继迁去世，李德明继承部众。此项以《宋史》卷485〈夏国传〉；《辽史》《金史》相关本纪；黑水城西夏文文书与考古报告 为基本来源线索；编年与官修材料能够确认年序和制度用语，但对地方执行、人数、动机或社会后果的判断仍须同文书、碑刻、考古及对方叙事互校。

\paragraph{1004年}\eventanchor{XIA-1004-REGNAL-YEAR}{XIA-1004-REGNAL-YEAR}　党项李氏以李继迁在位第23年作为诏令、赋役与外交文书的纪年框架。此项以《宋史》卷485〈夏国传〉；《辽史》《金史》相关本纪；黑水城西夏文文书与考古报告 为基本来源线索；编年与官修材料能够确认年序和制度用语，但对地方执行、人数、动机或社会后果的判断仍须同文书、碑刻、考古及对方叙事互校。

\paragraph{1005年}\eventanchor{XIA-1005-REGNAL-YEAR}{XIA-1005-REGNAL-YEAR}　党项李氏以李德明在位第1年作为诏令、赋役与外交文书的纪年框架。此项以《宋史》卷485〈夏国传〉；《辽史》《金史》相关本纪；黑水城西夏文文书与考古报告 为基本来源线索；编年与官修材料能够确认年序和制度用语，但对地方执行、人数、动机或社会后果的判断仍须同文书、碑刻、考古及对方叙事互校。

\paragraph{1006年}\eventanchor{XIA-1006-REGNAL-YEAR}{XIA-1006-REGNAL-YEAR}　党项李氏以李德明在位第2年作为诏令、赋役与外交文书的纪年框架。此项以《宋史》卷485〈夏国传〉；《辽史》《金史》相关本纪；黑水城西夏文文书与考古报告 为基本来源线索；编年与官修材料能够确认年序和制度用语，但对地方执行、人数、动机或社会后果的判断仍须同文书、碑刻、考古及对方叙事互校。

\paragraph{1007年}\eventanchor{XIA-1007-REGNAL-YEAR}{XIA-1007-REGNAL-YEAR}　党项李氏以李德明在位第3年作为诏令、赋役与外交文书的纪年框架。此项以《宋史》卷485〈夏国传〉；《辽史》《金史》相关本纪；黑水城西夏文文书与考古报告 为基本来源线索；编年与官修材料能够确认年序和制度用语，但对地方执行、人数、动机或社会后果的判断仍须同文书、碑刻、考古及对方叙事互校。

\paragraph{1008年}\eventanchor{XIA-1008-REGNAL-YEAR}{XIA-1008-REGNAL-YEAR}　党项李氏以李德明在位第4年作为诏令、赋役与外交文书的纪年框架。此项以《宋史》卷485〈夏国传〉；《辽史》《金史》相关本纪；黑水城西夏文文书与考古报告 为基本来源线索；编年与官修材料能够确认年序和制度用语，但对地方执行、人数、动机或社会后果的判断仍须同文书、碑刻、考古及对方叙事互校。

\paragraph{1009年}\eventanchor{XIA-1009-EVENT-052}{XIA-1009-EVENT-052}　李德明接受宋廷名号并经营夏州、灵州一带。此项以《宋史》卷485〈夏国传〉；《辽史》《金史》相关本纪；黑水城西夏文文书与考古报告 为基本来源线索；编年与官修材料能够确认年序和制度用语，但对地方执行、人数、动机或社会后果的判断仍须同文书、碑刻、考古及对方叙事互校。

\paragraph{1009年}\eventanchor{XIA-1009-REGNAL-YEAR}{XIA-1009-REGNAL-YEAR}　党项李氏以李德明在位第5年作为诏令、赋役与外交文书的纪年框架。此项以《宋史》卷485〈夏国传〉；《辽史》《金史》相关本纪；黑水城西夏文文书与考古报告 为基本来源线索；编年与官修材料能够确认年序和制度用语，但对地方执行、人数、动机或社会后果的判断仍须同文书、碑刻、考古及对方叙事互校。

\paragraph{1010年}\eventanchor{XIA-1010-REGNAL-YEAR}{XIA-1010-REGNAL-YEAR}　党项李氏以李德明在位第6年作为诏令、赋役与外交文书的纪年框架。此项以《宋史》卷485〈夏国传〉；《辽史》《金史》相关本纪；黑水城西夏文文书与考古报告 为基本来源线索；编年与官修材料能够确认年序和制度用语，但对地方执行、人数、动机或社会后果的判断仍须同文书、碑刻、考古及对方叙事互校。

\paragraph{1011年}\eventanchor{XIA-1011-REGNAL-YEAR}{XIA-1011-REGNAL-YEAR}　党项李氏以李德明在位第7年作为诏令、赋役与外交文书的纪年框架。此项以《宋史》卷485〈夏国传〉；《辽史》《金史》相关本纪；黑水城西夏文文书与考古报告 为基本来源线索；编年与官修材料能够确认年序和制度用语，但对地方执行、人数、动机或社会后果的判断仍须同文书、碑刻、考古及对方叙事互校。

\paragraph{1012年}\eventanchor{XIA-1012-REGNAL-YEAR}{XIA-1012-REGNAL-YEAR}　党项李氏以李德明在位第8年作为诏令、赋役与外交文书的纪年框架。此项以《宋史》卷485〈夏国传〉；《辽史》《金史》相关本纪；黑水城西夏文文书与考古报告 为基本来源线索；编年与官修材料能够确认年序和制度用语，但对地方执行、人数、动机或社会后果的判断仍须同文书、碑刻、考古及对方叙事互校。

